% =====================================================================
% APPENDIX A: USE OF AI SCIENTIFIC ASSISTANTS
% =====================================================================
\appendixbanner{APPENDIX A: USE OF AI SCIENTIFIC ASSISTANTS}
\addcontentsline{toc}{section}{Appendix A: Use of AI scientific assistants}

\subsection*{Workflow}

Three AI assistants were used for the theoretical part, each with a defined role and a verification
step (the full log is \texttt{docs/\allowbreak{}AI\_\allowbreak{}USAGE\_\allowbreak{}LOG.\allowbreak{}md} in the repository):
\begin{enumerate}
  \item \textbf{Perplexity} searched the literature with structured prompts that asked, for every
        source, the full reference, the DOI, open access and the exact claim it supports.
  \item \textbf{NotebookLM}, restricted to the downloaded PDFs, produced (a) an inventory of every file
        with the bibliographic data as printed in the document, and (b) for each topic of the report,
        short statements with the file, the page and a verbatim quote of at most 30 words.
  \item \textbf{Claude} decided which sources to keep from the inventory, wrote the text using only the
        quote-supported part of each statement, marked every claim still lacking a quote with a
        visible \texttt{[verify]} tag, and drafted the interpretation of the results, which the author
        checked against the notebooks.
\end{enumerate}
The author checks each tagged claim against the PDF before the final version.

\subsection*{Errors detected (candidates for the AI Error Bounty)}

\begin{enumerate}
  \item \textbf{Off-topic sources.} For a request on glioma classification, Perplexity returned a
        cardiac epigenetic reprogramming paper (twice) and a study on lipid metabolism in glioblastoma
        cell lines, neither supporting any requested point. Detected by comparing title and topic with
        the request; removed.
  \item \textbf{Wrong title and authors, repeated by two tools.} Perplexity and then NotebookLM gave
        one of the downloaded files as ``Machine Learning and Statistical Evaluation of Glioma
        Datasets'' by Garbulowski and Komorowski; the extraction step later attributed its content to
        Sánchez-Marqués et al., and the DOAJ record confirms the latter title and authors
        \cite{sanchezmarques2024}. Two independent AI tools agreeing is therefore not a verification.
  \item \textbf{Misattribution.} Parsons et al.\ (2008, \emph{Science}) was presented as a TCGA
        publication by both tools; it was removed from the TCGA sources.
  \item \textbf{Outdated notion of time.} NotebookLM flagged a 2025 publication year as ``a future
        publication year'', although the inventory was made in October 2026.
  \item \textbf{Wrong DOI.} NotebookLM gave 10.1093/neuonc/noad108 for the EANO guideline on molecular
        diagnostic tools; a Crossref check of all DOIs (script \texttt{report/\allowbreak{}check\_\allowbreak{}references.\allowbreak{}py})
        showed that it belongs to an unrelated article on CAR-T cells, and the correct DOI
        (noad100) was confirmed on the EANO website \cite{sahm2023}. The same check corrected the title
        that NotebookLM gave for \cite{varma2006}.
  \item \textbf{One paper, two identities.} Two uploaded files shared the same DOI; the inventory gave
        one of them a different title and garbled authors, so it was discarded as unreliable, while
        the other was correctly identified later \cite{leiriademendonca2025}.
  \item \textbf{Statements broader than their quotes.} Several extracted statements contained numbers
        or claims absent from the supporting quote (for example, glioblastoma incidence and survival
        figures attributed to a file whose quote only gave lower-grade glioma survival). Only the
        quoted part was used. Two claims of this kind were resolved with a second, targeted query
        that asked for the exact quote: an extracted statement that one study had not tuned its
        models inside the cross-validation was supported only by a quote on feature selection, and the
        targeted query showed that the classifiers used default hyperparameters \cite{tasci2022}; a
        count of reclassified cases, presented as a change from the original TCGA labels, compared
        two re-annotation pipelines instead \cite{leiriademendonca2025}.
  \item \textbf{A plausible but wrong interpretation (Claude).} When drafting the regression results,
        the assistant read the two groups of out-of-fold predictions as the grade (``the models recover
        the grade and predict the typical age of each grade''), supported only by a correlation of
        \RegCorrGrade{} with the grade, and wrote the conclusion of H2 accordingly. \emph{Detection}: the
        causal analysis estimated that only \CsPropPct\% of the IDH--age association goes through the
        grade, which contradicted that reading; checking the alternative explanation showed a stronger
        correlation with the IDH status (\RegCorrIDH) and the same split of predictions within each
        grade (Table~\ref{tab:prdl-reg-idh}). \emph{Correction}: the text, the conclusion of H2 and the
        regression notebook (which now reports both correlations) were rewritten. A correlation between
        a prediction and a variable that was not an input does not show that the model uses that
        variable when it is correlated with another one.
\end{enumerate}

\subsection*{AI Error Bounty sheets}

The strongest errors are documented below in the format of the course guideline. They were found in
the verification loop described above, in which every output of the assistants was checked against its
source, the data or a second run; each check is reproduced by the author from the steps in
\texttt{docs/\allowbreak{}BOUNTY\_\allowbreak{}CHECKS.\allowbreak{}md}.

\subsubsection*{Sheet P1. A wrong DOI pointing to another paper (NotebookLM)}
\begin{description}[style=nextline, leftmargin=1.2em]
  \item[Prompt] Inventory of the uploaded PDFs: one row per file with the bibliographic data ``as printed
        in the document'', writing ``not stated'' instead of guessing.
  \item[Faulty response] For the EANO guideline on molecular diagnostic tools (Sahm et al., 2023) the
        DOI 10.1093/neuonc/noad108.
  \item[Proof] The script \texttt{report/\allowbreak{}check\_\allowbreak{}references.\allowbreak{}py} compares the title and year of every DOI
        with the bibliography. The DOI resolved to an unrelated article on CAR-T cells in
        glioblastoma; the EANO page gives 10.1093/neuonc/noad100.
  \item[Correction] DOI replaced \cite{sahm2023}; every DOI of the bibliography is now checked with the
        script.
  \item[Reproduction] \todoverify{author: open both DOIs and record the titles}
  \item[Reflection] A DOI looks like hard data, but here it was generated, not read: the metadata of a
        source need the same verification as its claims.
\end{description}

\subsubsection*{Sheet P2. The same wrong authors from two tools (Perplexity and NotebookLM)}
\begin{description}[style=nextline, leftmargin=1.2em]
  \item[Prompt] Perplexity: machine-learning studies on glioma grading with clinical and mutation data,
        with full reference and DOI. NotebookLM: the inventory prompt of Sheet P1.
  \item[Faulty response] Both tools gave one of the downloaded papers as ``Machine Learning and
        Statistical Evaluation of Glioma Datasets'' by Garbulowski and Komorowski.
  \item[Proof] The claim-extraction step of NotebookLM attributed the content of the same file to
        Sánchez-Marqués et al.; the DOAJ record of its DOI confirms that title and authors. The two outputs of the same tool contradicted each other.
  \item[Correction] Reference corrected \cite{sanchezmarques2024}; agreement between AI tools is no
        longer treated as verification.
  \item[Reproduction] \todoverify{author: read the first page of 21.pdf}
  \item[Reflection] Two assistants agreeing is not two independent sources: the second may simply have
        repeated the first.
\end{description}

\subsubsection*{Sheet P3. Statements broader than their supporting quotes (NotebookLM)}
\begin{description}[style=nextline, leftmargin=1.2em]
  \item[Prompt] For each topic of the report: short statements, each with file, page and a verbatim
        quote of at most 30 words.
  \item[Faulty response] (a) ``Neither paper performed nested CV'' for Tasci et al., supported by a quote
        about feature selection; (b) ``811 samples maintained consistent classification, 134 changed''
        presented as the change from the original TCGA labels to the 2021 classification.
  \item[Proof] Each statement was compared with its own quote, and the quotes did not contain the
        claims. A second, targeted query showed that (a) Tasci et al.\
        used ``default values'' for the hyperparameters \cite{tasci2022}, and (b) the two compared
        classifications were two re-annotation pipelines, neither of them the original label
        \cite{leiriademendonca2025}.
  \item[Correction] Only the quoted content is used in Sections~2 and~6, with the caveats.
  \item[Reproduction] \todoverify{author: read 19.pdf p.~5 and 42.pdf p.~3}
  \item[Reflection] Asking for a quote is useful only if the quote is then read against the statement;
        the summary around a correct quote can still be wrong.
\end{description}

\subsubsection*{Sheet P4. A plausible but wrong interpretation of a figure (Claude)}
\begin{description}[style=nextline, leftmargin=1.2em]
  \item[Prompt] Regression notebook with the out-of-fold predictions and their interpretation; later,
        drafting of the results section of this report.
  \item[Faulty response] ``The models recover the grade through IDH1, IDH2 and ATRX and predict the
        typical age of each grade'', based on a correlation of \RegCorrGrade{} between the predictions
        and the grade; the conclusion of H2 said that the mutations explain age ``largely through the
        grade''.
  \item[Proof] The causal analysis estimated that only \CsPropPct\% of the
        IDH--age association goes through the grade. Testing the alternative explanation:
        the predictions correlate \RegCorrIDH{} with the IDH status, and within GBM the IDH-mutant
        patients are predicted at \RegPredGBMIDHmut{} years against \RegPredGBMIDHwt{} for IDH-wildtype
        (Table~\ref{tab:prdl-reg-idh}).
  \item[Correction] Notebook 04 reports both correlations and the grade $\times$ IDH table; the
        discussion and the conclusion of H2 were rewritten.
  \item[Reproduction] \todoverify{author: run notebook 04, section 6, and compare}
  \item[Reflection] A correlation between a prediction and a variable that was not an input does not
        show that the model uses that variable when another variable is correlated with both.
\end{description}

\subsection*{Reflection}

The verification loop changed how the sources were used: the bibliography records, for every entry,
how its metadata were checked, and no claim enters the text on the authority of an assistant alone.
